\chapter{Divider (fdiv)}
\label{ch:fdiv}

\section{Introduction}

\sig{fdiv} computes $\mathtt{op\_a}/\mathtt{op\_b}$. Two \sig{unpack} modules
split the operands. The \sig{div} module normalizes the mantissas and divides them
with an iterative restoring divider, which produces $\mathtt{MAN\_W}+3$ quotient
bits, one per clock cycle, and a sticky bit from the remainder. \sig{norm}
normalizes the quotient, and \sig{pack} rounds it and forms the result, as
Figure~\ref{fig:fdiv} shows.

\begin{figure}[H]
\centering
\begin{tikzpicture}[node distance=7mm]
  \node[ugblock] (u) {\sig{unpack}\\(2)};
  \node[ugblock,right=of u,minimum width=30mm] (a) {\sig{div}};
  \node[ugblock,right=of a] (n) {\sig{norm}};
  \node[ugblock,right=of n] (p) {\sig{pack}};
  \draw[ugarrow] ($(u.west)-(9mm,0)$) node[above right,font=\footnotesize\ttfamily,xshift=-1mm] {op} -- (u);
  \draw[ugarrow] (u) -- (a); \draw[ugarrow] (a) -- (n); \draw[ugarrow] (n) -- (p);
  \draw[ugarrow] (p) -- ($(p.east)+(9mm,0)$) node[above left,font=\footnotesize\ttfamily,xshift=1mm] {res};
\end{tikzpicture}
\caption{\sig{fdiv} Block Diagram}
\label{fig:fdiv}
\end{figure}

\section{Features}

\begin{itemize}
\item Iterative division with a latency of $\mathtt{MAN\_W}+6$ clock
  cycles: 30 for binary32.
\item One operation at a time.
\item Subnormal operands and results.
\item Round to nearest, ties to even.
\item NaN, infinity and zero-divisor handling, selected by \sig{SPECIAL\_CASES}.
\end{itemize}

\section{Parameters}

\sig{fdiv} has the parameters of Table~\ref{tab:common_params}.

\section{Interface Signals}

\sig{fdiv} has the signals of Table~\ref{tab:common_sigs}.

\Needspace*{23\baselineskip}
\section{Instantiation}

The example below instantiates \sig{fdiv} for binary32. Compile it with the files
of Table~\ref{tab:files}.

\begin{vcode}
fdiv #(
   .DATA_W       (32),
   .EXP_W        (8),
   .SPECIAL_CASES(1)
) fdiv0 (
   .clk    (clk),
   .rst    (rst),
   .valid_i(valid_i),
   .valid_o(valid_o),
   .op_a   (op_a),
   .op_b   (op_b),
   .res    (res)
);
\end{vcode}

\section{Operation}

\sig{fdiv} follows the protocol of Section~\ref{sec:protocol}, with these
differences:

\begin{itemize}
\item \textbf{Operands.} The unit samples \sig{op\_a} and \sig{op\_b} in the
  \sig{valid\_i} cycle, so they can change in the next cycle.
\item \textbf{Latency.} The latency is $\mathtt{MAN\_W}+6$ clock cycles,
  counted from the last cycle in which \sig{valid\_i} is high. It was measured as
  17, 30 and 59 cycles for the 16-, 32- and 64-bit formats ($\mathtt{MAN\_W}$ of
  11, 24 and 53), for normal, subnormal, zero, infinite and NaN operands.
\item \textbf{Next operation.} A new division can start in the clock cycle in
  which \sig{valid\_o} is high.
\item \textbf{Restart.} A \sig{valid\_i} during a division abandons it: only the new
  division produces a \sig{valid\_o} pulse.
\end{itemize}

Figure~\ref{fig:timing_div} shows a binary32 division.

\begin{figure}[H]
\centering
\begin{tikztimingtable}[timing/.style={x=3.2mm,y=3mm},timing/d/text/.style={font=\scriptsize}]
cycle & 2D{0} 2D{1} 6D{$\cdots$} 2D{29} 2D{30} 2D{31} \\
\sig{clk} & 2{HL} 6D{$\cdots$} 3{HL} \\
\sig{valid\_i} & 2H 14L \\
\sig{op\_a}, \sig{op\_b} & 2D{op} 14U \\
\sig{valid\_o} & 12L 2H 2L \\
\sig{res} & 12U 2D{res} 2U \\
\end{tikztimingtable}
\caption{Timing Diagram, Binary32 Division}
\label{fig:timing_div}
\end{figure}

Table~\ref{tab:fdiv_special} lists the results for special operands when
\sig{SPECIAL\_CASES} is 1. The sign of an infinite or zero result is the
exclusive OR of the operand signs.

\begin{twotable}{\sig{fdiv} Special Results}{tab:fdiv_special}{Operands}{Result}{7cm}
Either operand is NaN & NaN \\ \hline
Infinity / infinity & NaN \\ \hline
Zero / zero & NaN \\ \hline
Finite number / infinity & Zero \\ \hline
Nonzero number or infinity / zero & Infinity \\ \hline
Infinity / finite nonzero number & Infinity \\
\end{twotable}
